%%============================================================================
%% REGULAMENTO ELEITORAL — Grupo de Usuários Linux e Software Livre
%%                         Estatística e Ciência de Dados - UFPR
%%============================================================================

\documentclass[
  12pt,          % tamanho base da fonte
  a4paper,       % formato de papel
  oneside,       % impressão frente única
  brazil         % idioma principal: português do Brasil
]{article}

%---------------------------------------------------------------------------
% Pacotes essenciais
%---------------------------------------------------------------------------
\usepackage[T1]{fontenc}            % codificação de fonte
\usepackage[utf8]{inputenc}         % entrada UTF-8 (pdflatex)
\usepackage{lmodern}                % fontes Latin Modern (vetoriais)
\usepackage{microtype}              % microtipografia (melhora legibilidade)

\usepackage[brazil]{babel}          % hifenização e nomes de seções em PT-BR

\usepackage[
  margin=2.5cm,
  top=3cm,
  bottom=3cm,
  headheight=25pt
]{geometry}                         % margens do documento

\usepackage{setspace}               % controle de espaçamento entre linhas
\usepackage{xcolor}                 % suporte a cores
\usepackage{titlesec}               % formatação de títulos de seções
\usepackage{titling}                % formatação do título do documento
\usepackage{fancyhdr}               % cabeçalhos e rodapés personalizados
\usepackage{enumitem}               % personalização de listas
\usepackage{lettrine}               % capitulares (letra inicial decorada)
\usepackage{graphicx}               % suporte a imagens (se necessário)
\usepackage{hyperref}               % hiperlinks e metadados do PDF
\usepackage{bookmark}               % bookmarks do PDF mais limpos
\usepackage{listings}               % listagem de código

%---------------------------------------------------------------------------
% Cores institucionais
%---------------------------------------------------------------------------
\definecolor{gulecdBlue}{HTML}{1B3A5C}   % azul escuro institucional
\definecolor{gulecdAccent}{HTML}{2E86AB} % azul de destaque
\definecolor{gulecdGray}{HTML}{555555}   % cinza para texto secundário
\definecolor{gulecdLight}{HTML}{F5F7FA}  % fundo claro para caixas
\definecolor{gulecdRule}{HTML}{D0D5DD}   % cor de filetes e réguas

%---------------------------------------------------------------------------
% Metadados do PDF
%---------------------------------------------------------------------------
\hypersetup{
  pdftitle    = {Regulamento Eleitoral — GULECD/UFPR},
  pdfauthor   = {Grupo de Usuários Linux em Estatística e Ciência de Dados - UFPR},
  pdfsubject  = {Regulamento Eleitoral institucional do GULECD},
  pdfkeywords = {regulamento eleitoral, software livre, código aberto, ciência de dados, UFPR},
  colorlinks  = true,
  linkcolor   = gulecdAccent,
  urlcolor    = gulecdAccent,
  citecolor   = gulecdAccent
}

%---------------------------------------------------------------------------
% Configurações tipográficas
%---------------------------------------------------------------------------
\onehalfspacing                     % espaçamento entre linhas = 1,5
\setlength{\parindent}{1.5cm}       % recuo de parágrafo
\setlength{\parskip}{0.5em}         % espaço entre parágrafos

%---------------------------------------------------------------------------
% Formato do título do documento
%---------------------------------------------------------------------------
\pretitle{%
  \begin{center}
  \vspace*{1cm}
  \rule{\textwidth}{2pt}\vspace{0.3cm}\\
  \LARGE\bfseries\color{gulecdBlue}
}
\posttitle{%
  \\\vspace{0.3cm}
  \rule{\textwidth}{2pt}
  \end{center}
  \vspace{0.5cm}
}
\preauthor{\begin{center}\large\color{gulecdGray}}
\postauthor{\end{center}}
\predate{\begin{center}\small\color{gulecdGray}}
\postdate{\end{center}}

%---------------------------------------------------------------------------
% Formato das seções
%---------------------------------------------------------------------------
\titleformat{\section}
  [block]
  {\normalfont\Large\bfseries\color{gulecdBlue}}
  {\thesection.}{0.8em}{}
  [\vspace{0.2em}{\color{gulecdRule}\rule{\linewidth}{0.6pt}}]

\titleformat{\subsection}
  [block]
  {\normalfont\large\bfseries\color{gulecdAccent}}
  {\thesubsection.}{0.8em}{}

\titlespacing*{\section}
  {0pt}{2.5ex plus 1ex minus .2ex}{1.5ex plus .2ex}
\titlespacing*{\subsection}
  {0pt}{2ex plus 1ex minus .2ex}{1ex plus .2ex}

%---------------------------------------------------------------------------
% Cabeçalho e rodapé
%---------------------------------------------------------------------------
\pagestyle{fancy}
\fancyhf{}
\fancyhead[L]{\small\color{gulecdGray}\textsc{GULECD/UFPR}}
\fancyhead[R]{\small\color{gulecdGray}\textit{Regulamento Eleitoral}}
\fancyfoot[C]{\small\color{gulecdGray}---\ \thepage\ ---}
\renewcommand{\headrulewidth}{0.4pt}
\renewcommand{\headrule}{%
  \hbox to\headwidth{%
    \color{gulecdRule}\leaders\hrule height \headrulewidth\hfill}}

% Remove cabeçalho/rodapé da primeira página (será redefinido manualmente)
\fancypagestyle{plain}{%
  \fancyhf{}%
  \fancyfoot[C]{\small\color{gulecdGray}---\ \thepage\ ---}%
  \renewcommand{\headrulewidth}{0pt}%
}

%---------------------------------------------------------------------------
% Aumentar profundidade máxima de listas
%---------------------------------------------------------------------------
\setlistdepth{6}

%---------------------------------------------------------------------------
% Ambiente personalizado para listas
%---------------------------------------------------------------------------
\newlist{conduta}{itemize}{4}
\setlist[conduta,1]{
  label      = {\color{gulecdAccent}\textbullet},
  leftmargin = 1.5em,
  itemsep    = 0.3em,
  topsep     = 0.4em
}
\setlist[conduta,2]{
  label      = {\color{gulecdAccent}\textendash},
  leftmargin = 1.2em,
  itemsep    = 0.2em,
  topsep     = 0.3em
}
\setlist[conduta,3]{
  label      = {\color{gulecdAccent}\cdot},
  leftmargin = 1em,
  itemsep    = 0.2em,
  topsep     = 0.3em
}
\setlist[conduta,4]{
  label      = {\color{gulecdAccent}-},
  leftmargin = 0.8em,
  itemsep    = 0.15em,
  topsep     = 0.2em
}

%---------------------------------------------------------------------------
% Ambiente para listas numeradas
%---------------------------------------------------------------------------
\newlist{numerado}{enumerate}{3}
\setlist[numerado,1]{
  label      = \arabic*.,
  leftmargin = 1.5em,
  itemsep    = 0.3em,
  topsep     = 0.4em
}
\setlist[numerado,2]{
  label      = \arabic*.,
  leftmargin = 1.2em,
  itemsep    = 0.2em,
  topsep     = 0.3em
}
\setlist[numerado,3]{
  label      = \arabic*.,
  leftmargin = 1em,
  itemsep    = 0.15em,
  topsep     = 0.2em
}

%---------------------------------------------------------------------------
% Configuração do listings para código
%---------------------------------------------------------------------------
\lstset{
  basicstyle=\ttfamily\small,
  backgroundcolor=\color{gulecdLight},
  frame=single,
  framerule=0.5pt,
  rulecolor=\color{gulecdRule},
  breaklines=true,
  breakatwhitespace=true,
  showstringspaces=false,
  tabsize=2,
  language=bash
}

%---------------------------------------------------------------------------
% Informações do documento
%---------------------------------------------------------------------------
\title{Regulamento Eleitoral\\[0.3em]
  \Large\mdseries\color{gulecdGray}%
  Grupo de Usuários Linux e Software Livre\\Estatística e Ciência de Dados - UFPR}
\date{\today}

%============================================================================
% INÍCIO DO DOCUMENTO
%============================================================================
\begin{document}

%--- Capa / Título principal -------------------------------------------------
\thispagestyle{empty}
\maketitle

\newpage

%--- Seção 1: Disposições Gerais --------------------------------------------
\section{Disposições Gerais}

\lettrine[lines=3, lhang=0.1, nindent=0em, slope=-0.5em]{E}{ste} regulamento estabelece as bases formais para a condução do processo eleitoral do GULECD/UFPR,
 garantindo que a escolha de sua liderança ocorra de maneira transparente, participativa e alinhada aos princípios de organização horizontal e rigor institucional,
  assegurando legitimidade e continuidade organizacional.

O presente documento estabelece as normas, procedimentos e diretrizes para a eleição da Coordenação do GULECD/UFPR, composta por:

\begin{conduta}
	\item 1 Coordenador
	\item 1 Vice-Coordenador
	\item Núcleo de Coordenação
\end{conduta}

A eleição será realizada anualmente, por meio de chapas, assegurando transparência, participação ampla e legitimidade democrática.

%--- Seção 2: Estrutura da Chapa --------------------------------------------
\section{Estrutura da Chapa}

A estrutura em chapas tem como objetivo garantir coerência entre liderança e equipe de execução, incentivando que os membros escolham como liderança um conjunto
 organizado de pessoas com visão comum para o grupo. 

Cada chapa deverá obrigatoriamente conter:

\begin{conduta}
	\item 1 candidato a Coordenador
	\item 1 candidato a Vice-Coordenador
	\item Um conjunto de membros propostos para o Núcleo de Coordenação
\end{conduta}

\subsection{Núcleo de Coordenação}

\begin{conduta}
	\item Número livre de membros (recomendado: 3 a 8)
	\item Estrutura horizontal, sem hierarquia interna
	\item Função operacional e estratégica
\end{conduta}

%--- Seção 3: Atribuições dos Cargos ----------------------------------------
\section{Atribuições dos Cargos}

Busca-se preservar a horizontalidade do conjunto de liderança, concentrando a autoridade decisória apenas quando estritamente necessário.

\subsection{Coordenador}

\begin{conduta}
	\item Conduzir o direcionamento geral do grupo
	\item Atuar como ponto final de decisão em situações de impasse
	\item Representar institucionalmente o grupo
	\item Garantir alinhamento com o Manifesto e demais documentos
\end{conduta}

\subsection{Vice-Coordenador}

\begin{conduta}
	\item Apoiar o Coordenador em todas as funções
	\item Substituí-lo em sua ausência
	\item Assumir responsabilidades delegadas
	\item Auxiliar na articulação interna
\end{conduta}

\subsection{Núcleo de Coordenação}

\begin{conduta}
	\item Organizar e executar atividades do grupo
	\item Propor iniciativas e projetos
	\item Manter infraestrutura (GitHub, comunicação, etc.)
	\item Apoiar novos membros
	\item Participar de decisões estratégicas
	\item Atuar de forma colaborativa e horizontal
	
    Nota: as atribuições do Núcleo de Coordenação aplicam-se também ao Coordenador e ao Vice-Coordenador.
\end{conduta}

%--- Seção 4: Responsabilidades da Chapa ------------------------------------
\section{Responsabilidades da Chapa}

A candidatura em chapa implica responsabilidade coletiva sobre a condução do grupo, com apresentação de uma proposta coerente de atuação que demonstre 
capacidade organizacional, visão estratégica e compromisso com os princípios do grupo.

Cada chapa deverá apresentar:

\begin{conduta}
	\item Uma proposta de atuação (plano de gestão)
	\item Diretrizes para atividades do grupo durante o mandato
	\item Estratégias para crescimento, organização e impacto
	\item Compromisso com os princípios do Manifesto e Código de Conduta
\end{conduta}

%--- Seção 5: Elegibilidade -------------------------------------------------
\section{Elegibilidade}

Permitem-se candidaturas amplas, sem burocracia excessiva, mas incentiva-se que candidatos tenham histórico de envolvimento.

Podem se candidatar:

\begin{conduta}
	\item Membros ativos do grupo
	\item Participantes que demonstrem envolvimento e contribuição prévia
\end{conduta}

Não há restrição formal de tempo mínimo de participação no grupo.

%--- Seção 6: Comissão Eleitoral --------------------------------------------
\section{Comissão Eleitoral}

A condução do processo eleitoral será responsabilidade da gestão vigente, garantindo continuidade institucional. 
Ao mesmo tempo, essa condução deve ser pautada por neutralidade e transparência.

A Coordenação atualmente em exercício atuará como comissão eleitoral responsável por:

\begin{conduta}
	\item Definir e divulgar o cronograma
	\item Validar candidaturas
	\item Garantir transparência do processo
	\item Organizar a votação
	\item Apurar e divulgar resultados
\end{conduta}

%--- Seção 7: Cronograma Eleitoral ------------------------------------------
\section{Cronograma Eleitoral}

O processo eleitoral deverá seguir o seguinte cronograma mínimo:

\subsection{Abertura do Processo}

\begin{conduta}
	\item Divulgação oficial da eleição
	\item Prazo máximo: 30 dias antes da data de encerramento da gestão
\end{conduta}

\subsection{Período de Inscrição de Chapas}

\begin{conduta}
	\item Duração: 7 dias
	\item Submissão via:
	\begin{conduta}
		\item GitHub (issue no repositório apropriado), ou
		\item Mensagem formal no grupo e envio dos dados necessários à atual coordenação
	\end{conduta}
\end{conduta}

A inscrição deve conter:

\begin{conduta}
	\item Nome da chapa
	\item Composição completa
	\item Plano de gestão
\end{conduta}

\subsection{Validação das Chapas}

\begin{conduta}
	\item Duração: até 2 dias
	\item Verificação de conformidade
\end{conduta}

\subsection{Período de Campanha}

\begin{conduta}
	\item Duração: 7 dias
	\item Apresentação de propostas
	\item Discussões abertas no grupo
\end{conduta}

\subsection{Votação}

\begin{conduta}
	\item Duração: 2 dias
	\item Realizada no grupo oficial de WhatsApp
\end{conduta}

\subsection{Apuração e Resultado}

\begin{conduta}
	\item Divulgação em até 24 horas após encerramento
\end{conduta}

%--- Seção 8: Processo de Votação -------------------------------------------
\section{Processo de Votação}

O processo de votação deve ser simples, acessível e transparente, garantindo que todos os membros possam participar sem barreiras técnicas, ao mesmo tempo em 
que assegura clareza na apuração dos resultados.

\subsection{Forma de Votação}

A votação será realizada por votação direta no grupo de WhatsApp, podendo ocorrer via:

\begin{conduta}
	\item Enquete nativa do WhatsApp, ou
	\item Declaração explícita de voto (se necessário)
\end{conduta}

\subsection{Regras}

\begin{conduta}
	\item Cada membro tem direito a 1 voto
	\item Voto não é transferível
	\item Votação aberta e transparente
\end{conduta}

\subsection{Quórum}

\begin{conduta}
	\item A atual coordenação deverá assegurar que mais de 90\% dos membros do grupo foram informados da data de realização das eleições.
	\item Em caso de baixa participação (menos de 10\% dos membros do grupo como votantes) deve-se realizar nova votação.
\end{conduta}

\subsection{Critério de Vitória}

\begin{conduta}
	\item Vence a chapa com maior número de votos
	\item Em caso de empate:
	\begin{conduta}
		\item Nova votação entre as chapas empatadas
	\end{conduta}
\end{conduta}

%--- Seção 9: Transparência -------------------------------------------------
\section{Transparência}

A transparência é um princípio central do processo eleitoral. 
Todas as etapas devem ser documentadas e acessíveis aos membros, permitindo auditoria informal por qualquer participante.

Para garantir integridade do processo:

\begin{conduta}
	\item Todas as etapas devem ser publicamente documentadas
	\item As propostas devem ser acessíveis a todos os membros
	\item Resultados devem ser divulgados de forma clara
\end{conduta}

Deve-se utilizar o GitHub institucional do grupo para:

\begin{conduta}
	\item Registro das chapas
	\item Discussão de propostas
	\item Arquivamento histórico
\end{conduta}

%--- Seção 10: Mandato ------------------------------------------------------
\section{Mandato}

A definição de mandato temporal permite renovação periódica da liderança e adaptação contínua do grupo além de incentivar a participação ativa de novos membros.

\begin{conduta}
	\item Duração: 1 ano
	\item Início: imediatamente após a eleição
	\item Possibilidade de reeleição sem limitação formal
\end{conduta}

%--- Seção 11: Vacância e Substituição --------------------------------------
\section{Vacância e Substituição}

\begin{conduta}
	\item O Vice-Coordenador assume em caso de vacância do Coordenador
	\item O Núcleo pode ser parcialmente reconfigurado pela Coordenação
	\item Em caso de vacância do Coordenador e Vice-Coordenador o Núcleo de Coordenação pode eleger um novo Coodenador e Vice-Coordenador entre seus membros até o final da gestão da chapa.
	\item Em caso de vacância total, nova eleição deve ser convocada
\end{conduta}

%--- Seção 12: Registro de Chapas -------------------------------------------
\section{Registro de Chapas}

O processo de registro de chapas tem como objetivo garantir organização, transparência e rastreabilidade das candidaturas, 
utilizando a infraestrutura do GitHub como mecanismo formal de submissão, documentação e auditoria. 

A inscrição de chapas deverá ser realizada obrigatoriamente por meio do repositório oficial de eleições do grupo, mesmo se submetidas previamente
via grupo de Whatsapp. Nesse caso, a atual coordenação fica responsável por submeter a inscrição via Github.

\subsection{Repositório de Registro}

As chapas deverão utilizar o repositório:

\begin{conduta}
	\item \texttt{eleicoes-coordenacao}
\end{conduta}

Dentro deste repositório, cada processo eleitoral ocorrerá em um diretório específico, por exemplo:

\begin{conduta}
	\item \texttt{eleicao-gestao-01}
\end{conduta}

\subsection{Procedimento de Submissão}

Para efetivar a inscrição, a chapa deverá:

\begin{numerado}
	\item Clonar o repositório:
	\begin{lstlisting}
git clone https://github.com/<org>/eleicoes-coordenacao.git
    \end{lstlisting}
	\item Navegar até o diretório da eleição vigente:
	\begin{lstlisting}
eleicao-gestao-01/
    \end{lstlisting}
	\item Criar um novo diretório com o nome da chapa, no formato:
	\begin{lstlisting}
chapaXX
    \end{lstlisting}
	onde \texttt{XX} corresponde ao número identificador da chapa.
\end{numerado}

\subsection{Estrutura Obrigatória}

Dentro do diretório da chapa, deverão ser incluídos as seguintes informações e documentos:

\begin{conduta}
	\item \textbf{Plano de gestão}
	\item \textbf{Lista de membros}, contendo:
	\begin{conduta}
		\item Nome completo
		\item Cargo (Coordenador, Vice-Coordenador, Núcleo de Coordenação)
	\end{conduta}
	\item \textbf{Mini-biografia de cada membro}
	\item \textbf{Link para currículo Lattes de cada membro}
\end{conduta}

Todos os documentos devem:

\begin{conduta}
	\item Ser escritos em \textbf{LaTeX (.tex)}
	\item Ser acompanhados das versões compiladas em:
	\begin{conduta}
		\item PDF (.pdf)
		\item Markdown (.md)
		\item HTML (.html)
	\end{conduta}
\end{conduta}

\subsection{Submissão via Pull Request}

Após a organização dos arquivos, a chapa deverá:

\begin{numerado}
	\item Realizar commit das alterações
	\item Abrir um \textbf{Pull Request (PR)} para o repositório
\end{numerado}

A abertura do PR:

\begin{conduta}
	\item Constitui formalmente a \textbf{inscrição da chapa no processo eleitoral}
	\item Define a \textbf{data oficial de inscrição}
	\item Torna pública a candidatura e seus documentos
\end{conduta}

\subsection{Validação}

A comissão eleitoral deverá:

\begin{conduta}
	\item Verificar a conformidade estrutural da submissão
	\item Confirmar a presença de todos os documentos obrigatórios
	\item Validar a elegibilidade da chapa
\end{conduta}

Caso haja inconsistências, a chapa poderá ser notificada para correção dentro do prazo de inscrição.

\subsection{Transparência e Registro Permanente}

Todos os registros de chapas:

\begin{conduta}
	\item Permanecerão disponíveis no repositório
	\item Constituirão arquivo histórico das eleições
	\item Poderão ser consultados livremente pelos membros
\end{conduta}

%--- Seção 13: Disposições Finais -------------------------------------------
\section{Disposições Finais}

Este regulamento deve ser entendido como um instrumento vivo, passível de aprimoramento conforme a evolução do grupo. 
Sua aplicação deve sempre priorizar os princípios fundamentais de abertura, participação, eficiência e responsabilidade coletiva.

\begin{conduta}
	\item Este regulamento poderá ser revisado mediante discussão coletiva
	\item Situações não previstas serão resolvidas pela Coordenação vigente
\end{conduta}

%--- Espaço para assinaturas ------------------------------------------------
\vspace{2cm}
\begin{center}
	\rule{0.45\textwidth}{0.4pt}\hfill\rule{0.45\textwidth}{0.4pt}\\[0.3em]
	\small\color{gulecdGray}
	Assinatura dos membros fundadores
\end{center}

%============================================================================
\end{document}
%============================================================================
